\section{Experimental design}
\label{sec:experimental_design}

The experimental evaluation investigates planning performance across three distinct tiers: large-scale open-loop statistical simulation, closed-loop rolling replanning, and production software integration within Apollo Planning. Statistical experiments use deterministic scenario generators and fixed random seeds so that every method receives identical inputs, with collision and clearance evaluated against separately maintained ground-truth trajectories.

\subsection{Protocols and scenarios}
\label{sec:protocols_scenarios}

Each scenario uses a fixed horizon $T\in\{9,10,12,14\}$~s assigned by family: 9~s for crossing pedestrians and delayed crossings; 10~s for vehicle cut-ins, parked-with-oncoming traffic, and prediction noise; 12~s for narrow multi-obstacle situations; and 14~s for interleaved dynamic traffic. An episode times out when the clock reaches $T$ without achieving $s_{\mathrm{end}}\geq s_{\max}-1$~m. Four disjoint protocols are used.

\emph{Open-loop evaluation over 3,500 scenarios.} 500 seeds in each of seven families: crossing pedestrians, vehicle cut-ins, parked with oncoming traffic, narrow multi-obstacle, interleaved dynamic, prediction noise, delayed crossings. Parameters are sampled within family ranges. Planning and truth scenarios coincide except in prediction-noise, where bounded errors are injected only into planning input.

\emph{Fixed-section geometry over 4,000 instances.} 20 deterministic seeds with 200 fixed-section interval instances per seed; each has zero to eight possibly overlapping, nested, or road-crossing forbidden intervals. The scan is compared with exhaustive enumeration of all $2^k$ side assignments on center-band width.

\emph{Rolling-replanning evaluation over 700 scenarios.} 100 seeds per family. The ego vehicle is initialized from truth state $(s_0,l_0,v_0,a_0)$ at $t=0$ with a cold-start first cycle that does not use warm-started QP duals or committed homotopy. Every 0.5~s the first planned segment executes, obstacle states update, and the planner is called on shifted horizon $\max(2.0\text{ s},T-t)$ until the vehicle achieves arrival at $s\geq s_{\max}-1$~m or reaches $T$.

\emph{Finite-grid joint reference over 70 scenarios.} Stratified with 10 seeds per family and grid spacings $\Delta t=0.5$ s, $\Delta s=0.5$ m, $\Delta l=0.25$ m, $\Delta v=1.0$ m/s, beam width 20,000.

\subsection{Baselines and implementation}
\label{sec:baselines_impl}

The matched PVD baselines are defined by information exchanged at the constraint interface:
\begin{itemize}
  \item B0 is time-blind and uses the fixed local side policy.
  \item B1 adds arrival-time-dependent projection but retains obstacle-wise greedy side selection.
  \item B2 allows three arrival-time refinement iterations in the B1 policy.
  \item MIKU adds group-level continuous partition, Top-$K$ spatial homotopy, threat-aware margins, uncertainty-expanded occupancy, temporal homotopy, and at most two planning iterations with alternating updates.
\end{itemize}
All four share identical path and speed QP objectives, kinematic constraints, input, and metric implementation. Both piecewise-jerk QPs use OSQP at default tolerances with no per-call wall-clock cap on the 3,500-scenario runs. Baseline B2 uses the same QPs inside three arrival-time refinement iterations. Baseline B3 is a deterministic beam-search dynamic program on the $(t,s,l,v)$ grid with a beam width of 20,000. Its per-transition cost weights speed deviation from $v_0$, longitudinal acceleration, lateral offset, and lateral acceleration minus a progress reward; transitions that are kinematically inadmissible or collide with uncertainty-expanded obstacles are pruned before ranking. MIKU hyperparameters were fixed before evaluation and are reported completely in Table~\ref{tab:fixed_parameters}. Runs used an Intel i9-14900HX with 64 GB RAM. Apollo integration confirms compatibility with Apollo 11.0 Planning.

\begin{table}[t]
\centering
\caption{Fixed planning and uncertainty parameters used in all reported experiments.}
\label{tab:fixed_parameters}
\small
\begin{tabular}{lll}
\toprule
Parameter & Value & Role \\
\midrule
$(w_1,\ldots,w_5)$ & $(0.30,0.20,0.15,0.10,0.25)$ & threat-score weights \\
TTC breakpoints & $2.0,7.0$ s & urgency normalization \\
$(\delta_{\min},\delta_{\max})$ & $(0.10,0.40)$ m & adaptive obstacle buffer \\
$Q,K_s,\lambda_s$ & $3,3,0.05$ & spatial bands, Top-$K$, transition weight \\
$B$ & $8$ & temporal-graph beam width \\
$\eta_c$ & $0.10+0.05\min(|v_l|,2)$ s & occupancy-window expansion \\
$\epsilon_s$ & $0.05$ m & longitudinal ST guard \\
$(\epsilon_{s0},\epsilon_{l0})$ & $(0.35,0.18)$ m & initial prediction-error bounds \\
$(\epsilon_{vs},\epsilon_{vl})$ & $(0.30,0.12)$ m/s & prediction-error growth bounds \\
Path discretization & $0.5$ m & path-bound sampling \\
Speed discretization & $0.1$ s & speed-QP sampling \\
Planning iterations & at most $2$ & initial solve plus refinement \\
\bottomrule
\end{tabular}
\end{table}

\subsection{Outcomes and statistical analysis}
\label{sec:outcomes_stats}

The primary outcome is collision-free goal reaching: a run succeeds when collision-free over the horizon and reaching $s_{\mathrm{end}}\geq s_{\max}-1$ m. Collision rate uses truth rectangles; degraded-stop rate records runs that stop or fail to reach the goal without collision. Progress ratio, average speed, jerk RMS, minimum signed clearance, minimum longitudinal TTC, penalised travel time, and planning latency give safety, comfort, and computational context. Penalised travel time equals arrival time for collision-free arrivals and family-specific $T$ from Section~\ref{sec:protocols_scenarios} plus 5~s otherwise. Runtime includes the full method call; B2 includes all coordination iterations.

For each fixed seed we compute MIKU minus baseline differences before aggregation. Continuous metrics use 5,000-resample paired percentile bootstrap intervals; binary outcomes use exact McNemar tests; effect sizes are paired Cohen's $d_z$. Runtime is reported at the 50th, 95th, and 99th percentiles.
